
The non-Abelian anyon of the Kitaev spin liquid is a bound state of a $\mathbb{Z}_2$ flux and Majorana zero modes.  Here we have studied how these electrically neutral excitations in a Mott insulator respond to electric fields.

Microscopically, inhomogeneous spin correlations in Mott insulators may lead to charge fluctuations. In this scenario, these inhomogeneities are due to the $\mathbb{Z}_2$ fluxes which provide a source of the nearest neighboring spin correlation profile. To find the effective charge distribution we started from the three $t_{2g}$ orbital Hubbard-Kanamori model for Kitaev materials.  The effective spin-1/2 Hamiltonian is obtained by a canonical transformation into the low energy sector of one hole per site. Considering hopping beyond dominant, the effective spin Hamiltonian corresponds to the extended Kitaev model with dominant ferromagnetic Kitaev exchange and integrability broken Heisenberg and $\Gamma$ interactions. By this projection, the effective charge density to third-order perturbation is quadratic in the (next-) nearest neighboring spin correlation.

The extended Kitaev model is not solvable, and to find the correlations we applied a Majorana mean-field theory with local mean-field parameters. An extensive number of parameters are needed to consider inhomogeneous flux configurations. We found that this mean field is exact for the pure Kitaev model, and captures the qualitative physics for the Kitaev spin liquid phase for non-zero Heisenberg $J$, $\Gamma$, and magnetic field $\mb h$.  The mean-field approach allows us to go beyond perturbative couplings, however, we remain focused on the regions of not-too-large exchanges where the mean-field parameters for the Kitaev spin liquid phase are finite.

The results for the charge density are that it vanishes in a region with no vortices and it is finite around the vortex position. Importantly, vortices remain localized on sufficiently long time scales even in the presence of small perturbations around the Kitaev limit, as long as the system remains in the Kitaev spin liquid phase.
In our mean-field calculation, the vortex is found to be at fixed positions determined by the choice of the initial configuration on the self-consistent solution.


The charge density $|\delta n_l|$ on sites around the vortex is of the order of $10^{-6}$,  producing local electric fields near the detection limit of state-of-the-art atomic force microscopy \cite{Pereira2020,Wagner2019,Bian2021,Yin2021}. Importantly, here we use estimates for the hopping $t^{\alpha\beta}_{ij}$ and interaction parameters $U, J_H$ for bulk $\alpha$-RuCl$_3$, but the charge fluctuations can be greatly enhanced if the on-site repulsion $U$ is screened in a monolayer by the interaction with a substrate.


The charge density decays exponentially with the distance from the vortex. As long as the state remains in the Kitaev spin liquid phase, Heisenberg and $\Gamma$ exchanges do not qualitatively change the charge profile. However, we found an important dependence on the direction of the applied magnetic field. For a field perpendicular to the honeycomb plane, the charge profile is isotropic, while for fields with in-plane components, the charge profile is anisotropic. The direction of the anisotropy is perpendicular to the in-plane component of the magnetic field.

%The charge density decays exponentially with the distance from the vortex. As long as the state remains in the Kitaev spin liquid phase, Heisenberg and $\Gamma$ exchanges do not change qualitatively the charge profile. However, we found an important dependence on the applied magnetic field direction. For the field perpendicular to the honeycomb plane, the charge profile is isotropic, while for fields with in-plane components, the charge profile is anisotropic. The direction of the anisotropy is perpendicular to the in-plane component of the magnetic field.

We found that the charge profile of an isolated vortex can be characterized by the quadrupole tensor. %The anisotropy in the charge distribution is characterized by the difference between the in-plane eigenvalues of the tensor.
The difference between the in-plane eigenvalues of the tensor characterizes the anisotropy in the charge distribution.
Depending on the value of anisotropy, the interaction between the vortices, via the electrical quadrupole, can be attractive or repulsive. For low anisotropy, which is found for a magnetic field perpendicular to the plane, only repulsive interaction is found. By tilting the magnetic field along the $ac$ plane to the $\hat{\mb z}$ direction as considering high $\Gamma$ exchange the interaction is found to be attractive along some in-plane directions.
At long distances, the algebraic decay of the quadrupole-quadrupole interaction dominates over other sources of interaction between the vortices which are expected to decay exponentially in the non-Abelian phase.

On the matter of manipulating the Ising anyons of the Kitaev spin liquid phase with a local electric potential, we found that the gap to create a vortex is reduced in comparison to the vortex-free configuration. This allows one to attract a vortex in the extended Kitaev model with a STM tip in a controlled way.
By considering the Dzyaloshinskii-Moriya term, which is generically present since the local electric field breaks bond inversion symmetry, we found that one can close the vortex gap with a finite electric potential.


We only studied isolated vortices so far. It would be instructive to understand what happens when the vortices are brought close together.
In the absence of electric potential, the flux gap is greater in this configuration than in the isolated case.  It remains to be understood whether the electric potential is also able to close the vortex gap in this flux configuration.


Our research lays the foundation for the electrical detection and manipulation of Kitaev materials. Specifically, achieving precise control of Ising anyons within these materials would be a crucial step toward establishing a platform for topological quantum computation.
